.. \+\_\+intro\+:

The Hankel transform, also known as the Fourier–\+Bessel transform, is widely used in scattering. The transform is needed in many areas, e.\+g. to calculate scattering functions, to correct for multiple scattering, to simulate spin echo signals like those obtained on Larmor at the ISIS Neutron and Muon Source, as well as to qualitatively analyse phase contrast images (PCI) taken at the Diamond synchrotron source.

Small-\/angle scattering (SAS) is an experimental technique using X-\/rays (SAXS) or neutrons (SANS) to study the structure of materials on the nanoscale (roughly 1–1000 nm). SAS measures how radiation is scattered at very small angles, which tells us about large-\/scale structures inside the sample.

In this library, there are 14 different implementations of the Hankel transform, corresponding to different ways of approximating the Hankel transform integral.

+-\/-\/-\/-\/-\/-\/-\/-\/-\/-\/-\/-\/-\/-\/-\/-\/-\/-\/-\/-\/-\/-\/-\/---+-\/-\/-\/-\/-\/-\/-\/-\/-\/-\/-\/-\/-\/-\/-\/-\/-\/-\/---+-\/-\/-\/-\/-\/-\/-\/-\/-\/-\/-\/-\/-\/-\/-\/-\/-\/-\/-\/-\/-\/-\/-\/-\/-\/-\/-\/-\/-\/-\/-\/-\/-\/-\/-\/-\/---+ $\vert$ index (as in SASfit) $\vert$ Strategy Name $\vert$ Type $\vert$ +==========================+=====================+=======================================+ $\vert$ 0 $\vert$ DE\+\_\+\+Ooura $\vert$ Double-\/exponential quadrature $\vert$ +-\/-\/-\/-\/-\/-\/-\/-\/-\/-\/-\/-\/-\/-\/-\/-\/-\/-\/-\/-\/-\/-\/-\/---+-\/-\/-\/-\/-\/-\/-\/-\/-\/-\/-\/-\/-\/-\/-\/-\/-\/-\/---+-\/-\/-\/-\/-\/-\/-\/-\/-\/-\/-\/-\/-\/-\/-\/-\/-\/-\/-\/-\/-\/-\/-\/-\/-\/-\/-\/-\/-\/-\/-\/-\/-\/-\/-\/-\/---+ $\vert$ 1 $\vert$ DE\+\_\+\+Ogata $\vert$ Double-\/exponential quadrature $\vert$ +-\/-\/-\/-\/-\/-\/-\/-\/-\/-\/-\/-\/-\/-\/-\/-\/-\/-\/-\/-\/-\/-\/-\/---+-\/-\/-\/-\/-\/-\/-\/-\/-\/-\/-\/-\/-\/-\/-\/-\/-\/-\/---+-\/-\/-\/-\/-\/-\/-\/-\/-\/-\/-\/-\/-\/-\/-\/-\/-\/-\/-\/-\/-\/-\/-\/-\/-\/-\/-\/-\/-\/-\/-\/-\/-\/-\/-\/-\/---+ $\vert$ 6 to 11 $\vert$ DHT\+\_\+6 to DHT\+\_\+11 $\vert$ Digital filters $\vert$ +-\/-\/-\/-\/-\/-\/-\/-\/-\/-\/-\/-\/-\/-\/-\/-\/-\/-\/-\/-\/-\/-\/-\/---+-\/-\/-\/-\/-\/-\/-\/-\/-\/-\/-\/-\/-\/-\/-\/-\/-\/-\/---+-\/-\/-\/-\/-\/-\/-\/-\/-\/-\/-\/-\/-\/-\/-\/-\/-\/-\/-\/-\/-\/-\/-\/-\/-\/-\/-\/-\/-\/-\/-\/-\/-\/-\/-\/-\/---+ $\vert$ 12 $\vert$ QWE\+\_\+\+Key $\vert$ QWE with continued fraction expansion $\vert$ +-\/-\/-\/-\/-\/-\/-\/-\/-\/-\/-\/-\/-\/-\/-\/-\/-\/-\/-\/-\/-\/-\/-\/---+-\/-\/-\/-\/-\/-\/-\/-\/-\/-\/-\/-\/-\/-\/-\/-\/-\/-\/---+-\/-\/-\/-\/-\/-\/-\/-\/-\/-\/-\/-\/-\/-\/-\/-\/-\/-\/-\/-\/-\/-\/-\/-\/-\/-\/-\/-\/-\/-\/-\/-\/-\/-\/-\/-\/---+ $\vert$ 13 $\vert$ QWE\+\_\+\+Chave $\vert$ QWE with Shanks transformation $\vert$ +-\/-\/-\/-\/-\/-\/-\/-\/-\/-\/-\/-\/-\/-\/-\/-\/-\/-\/-\/-\/-\/-\/-\/---+-\/-\/-\/-\/-\/-\/-\/-\/-\/-\/-\/-\/-\/-\/-\/-\/-\/-\/---+-\/-\/-\/-\/-\/-\/-\/-\/-\/-\/-\/-\/-\/-\/-\/-\/-\/-\/-\/-\/-\/-\/-\/-\/-\/-\/-\/-\/-\/-\/-\/-\/-\/-\/-\/-\/---+ $\vert$ ... $\vert$ ... $\vert$ ... $\vert$ +-\/-\/-\/-\/-\/-\/-\/-\/-\/-\/-\/-\/-\/-\/-\/-\/-\/-\/-\/-\/-\/-\/-\/---+-\/-\/-\/-\/-\/-\/-\/-\/-\/-\/-\/-\/-\/-\/-\/-\/-\/-\/---+-\/-\/-\/-\/-\/-\/-\/-\/-\/-\/-\/-\/-\/-\/-\/-\/-\/-\/-\/-\/-\/-\/-\/-\/-\/-\/-\/-\/-\/-\/-\/-\/-\/-\/-\/-\/---+

The column named as \char`\"{}index\char`\"{} in the table essentially refers to the numbering used in SASfit.

Strategies 6 to 11 in SASfit, which in this library are referred to as DHT\+\_\+6, DHT\+\_\+7 .. DHT\+\_\+11, are fixed-\/grid digital filter strategies. They are quite fast and perform well on simple form factors but their accuracy decreases on more complex or oscillatory ones. They are not iterative, rely on pre-\/computed constants, and do not allow the user to control the \char`\"{}accuracy\char`\"{} of the solution in any way.

All other strategies instead are iterative, and some of them require the user to supply a tolerance parameter (called {\ttfamily eps\+\_\+rel}) which is then used to determine when to stop iterating. Specifically, iterations are stopped when the difference in the solution of two consecutive iterations is less than the tolerance. For this reasons, strategies that require eps\+\_\+rel generally perform better than the DHT ones on complex form factors, though they are more slow. Examples of these strategies are 1, 12 and 13 in the above table.

... 